\tzpanchor{0009}
\tzpentryheader{0009}{Hyper-Dimensional Span}

\begingroup\small
\begingroup\scriptsize
\setlength{\tabcolsep}{4pt}
\renewcommand{\arraystretch}{0.92}
\noindent\begin{tabularx}{\linewidth}{@{}p{0.24\linewidth}p{0.36\linewidth}X@{}}
\textbf{UUID} & \multicolumn{2}{l@{}}{\tzpcode{10d00b66-ffcf-5948-945f-34443de537a2}}\\
\textbf{PiID} & \tzpcode{5897932384626} & \textbf{TZPi}\quad \tzpcode{2}\quad \textbf{PiPE}\quad \tzpcode{16}\\
\textbf{RTEID} & \textbf{Poloidal $\theta$}\quad \tzpcode{34.9490 rad} & \textbf{Toroidal $\phi$}\quad \tzpcode{00.0576 rad}\\
\textbf{TZPID} & \tzpcode{ID0009} & \textbf{Node Dependencies}\quad \tzpidref{0003}\\
\textbf{Registry} & \tzpcode{registry\_id\_0009} & \textbf{Scale}\quad hyper-dimensional\\
\textbf{LOC Classification} & QA611.75 higher-dimensional geometry & QC174.17 field integration\\
\textbf{arXiv Classification} & Primary: hep-th \quad Cross-list: math-ph, gr-qc & \\
\textbf{Primary Support} & Branch Neighborhood Declaration & Hyper-Dimensional Measure\\
\textbf{Closure Support} & Normalized Span & \\
\end{tabularx}\par
\endgroup
\endgroup

\tzpsection{Canonical Statement}
The hyper-dimensional span measures the effective volume available in the neighborhood where local dimension is allowed to branch around the zero point. It is the geometric quantity that turns dimensional ambiguity into an integration domain for later actions and field measures.

\tzpsection{Canonical Equation}
\[
S_{\mathrm{HD}}
:=
\int_{\mathcal U}
\sqrt{\lvert \det h \rvert}\, d^{n}x
\]
\[
\mathcal U \ni p,
\qquad
\dim_{p}(M)\in\{n-1,n,n+1\}
\]

\begin{minipage}[t]{0.49\linewidth}
\tzpsection{Canonical Symbol Key}
\begingroup
\small
\raggedright
\textbf{Source equation symbols.}\par
\begin{itemize}[leftmargin=1.35em,itemsep=1pt,topsep=1pt,parsep=0pt]
    \item $S_{\mathrm{HD}}$: source equation symbol.
    \item $U$: auxiliary source state or local chart variable.
    \item $n$: integer mode index.
    \item $M$: source equation symbol.
    \item $u_{T}$: source equation symbol.
    \item $T$: source equation symbol.
\end{itemize}
\endgroup
\end{minipage}\hfill
\begin{minipage}[t]{0.47\linewidth}
\tzpsection{Wolfram Form}
\begingroup
\scriptsize
\raggedright
\textbf{Source Wolfram model.}\par
\begin{Verbatim}[fontsize=\tiny,breaklines=true,breakanywhere=true]
(* TZPID ID0009 generated source Wolfram model *)
ClearAll[K0009, Cr0009, T, R, A, W, I, N];
eq0009$1 = HoldForm[SHD == IntegralMathcalU*Sqrt[Abs[Det*h]]*d^n*x];
eq0009$2 = HoldForm[Element[Contains[mathcal*U, TZPSequence[p, DimP[M]]], {n - 1, n, n + 1}]];
eq0009$3 = HoldForm[SHDNorm == NormalizeTRAWIN[IntegralMathcalUSqrt[Abs[Det*h]]*d^n*x, DimP[M], uT]];

K0009 = HoldForm[{eq0009$1, eq0009$2, eq0009$3}];

N = "normalized TRAWIN closure object for Hyper-Dimensional Span";
I = "Normalized Span integration/comparison layer";
W = "Hyper-Dimensional Measure wave or operator carrier";
A = "Branch Neighborhood Declaration admissible action/amplitude condition";
R = "Hyper-Dimensional Measure rotation/curl preservation layer";
T = "Branch Neighborhood Declaration local source condition";

Cr0009[expr_] := HoldForm[N[I[W[A[R[T[expr]]]]]]];

Result0009 = Cr0009[K0009];
\end{Verbatim}
\endgroup
\end{minipage}

\par\vspace{0.45\baselineskip}
\noindent\begin{minipage}[t]{\linewidth}
\begingroup
\small
\raggedright
\textbf{TRAWIN Operator Key.}\par
\noindent\textbf{Time/localization operator:} $\mathsf{T}\!\left[\mathrm{local}\right]$ Branch Neighborhood Declaration local source condition.\par
\noindent\textbf{Rotation/curl operator:} $\mathsf{R}\!\left[\nabla\times\right]$ Hyper-Dimensional Measure rotation/curl preservation layer.\par
\noindent\textbf{Action/amplitude operator:} $\mathsf{A}\!\left[\mathrm{admissible}\right]$ Branch Neighborhood Declaration admissible action/amplitude condition.\par
\noindent\textbf{Wave/d'Alembertian operator:} $\mathsf{W}\!\left[\Box\right]$ Hyper-Dimensional Measure wave or operator carrier.\par
\noindent\textbf{Integration operator:} $\mathsf{I}\!\left[\int\right]$ Normalized Span integration/comparison layer.\par
\noindent\textbf{Normalization operator:} $\mathsf{N}\!\left[\mathrm{normalized}\right]$ normalized TRAWIN closure object for Hyper-Dimensional Span.\par
\endgroup
\end{minipage}

\clearpage
\noindent\textbf{[ID 0009] Hyper-Dimensional Span}\par
\vspace{0.35em}
\tzpsection{Derivation Layer}
\paragraph{Primary Equation.}
\begin{equation}
T_{\mathrm{IM}}
=
\frac{1}{u_T}
\int_{\Sigma}
\phi_{m_1}^{*}
\,\mathcal{L}\,
\phi_{m_2}
\,\mathrm{d}\Sigma .
\end{equation}

\paragraph{Primary Derivation.}
Let $\phi_{m_1}$ and $\phi_{m_2}$ be two modal states defined over a hypersurface $\Sigma$.  
Let $\mathcal{L}$ be an operator mediating transfer between these modes.

The raw modal transfer amplitude is given by the inner-product-like integral:

\begin{equation}
\int_{\Sigma}
\phi_{m_1}^{*}
\,\mathcal{L}\,
\phi_{m_2}
\,\mathrm{d}\Sigma .
\end{equation}

To express this transfer in Trawin-normalized units, divide by the base unit $u_T$:

\begin{equation}
T_{\mathrm{IM}}
=
\frac{1}{u_T}
\int_{\Sigma}
\phi_{m_1}^{*}
\,\mathcal{L}\,
\phi_{m_2}
\,\mathrm{d}\Sigma .
\end{equation}

\paragraph{Interpretation.}
This quantity measures inter-modal transfer across a hypersurface, scaled by the fundamental TZP loop unit.

\par\vspace{0.35em}
\noindent\textbf{Derivable Consequences.}\quad
\begingroup
\footnotesize
\raggedright
The canonical equation anchors Hyper-Dimensional Span by connecting the source condition to the derivation layer and proof certificate.
\par\endgroup

\tzpsection{Equation Support}
\noindent\textbf{1. Branch Neighborhood Declaration}\par
\[
\mathcal U \ni p,
\qquad
\dim_{p}(M)\in\{n-1,n,n+1\}
\]
\noindent This specifies the singular neighborhood where a higher-dimensional measure is even meaningful.\par
\noindent\textbf{2. Hyper-Dimensional Measure}\par
\[
S_{\mathrm{HD}}
:=
\int_{\mathcal U}
\sqrt{\lvert \det h \rvert}\, d^{n}x
\]
\noindent This converts the ambiguous local geometry into an explicit span that can be integrated and compared.\par
\noindent\textbf{3. Normalized Span}\par
\[
S_{\mathrm{HD}}^{\mathrm{norm}}
=
\operatorname{Normalize}_{\mathrm{TRAWIN}}
\!\left(
\int_{\mathcal U}\sqrt{\lvert \det h \rvert}\, d^{n}x,\,
\dim_{p}(M),\,
u_{T}
\right)
\]
\noindent This closes the progression by tying the span to both dimensional branching and the internal Trawin unit, making it transportable across admissible models.\par

\clearpage
\noindent\textbf{[ID 0009] Hyper-Dimensional Span}\par
\vspace{0.35em}
\tzpsection{Dictionary}
\begingroup\small
\tzpsection{Definition}
Hyper-Dimensional Span is a registry definition in the active TZPID arc.

% BEGIN TZP CONTEXTUAL MEANING
\tzpsection{Contextual Meaning}
\begingroup\footnotesize\setlength{\emergencystretch}{3em}\sloppy
\textbf{Formal statement:} T sub IM = ratio 1 u sub T integral sub Sigma phi sub m sub 1 to the power * L phi sub m sub 2 d Sigma,. \textbf{Contextual meaning:} The entry reads Hyper-Dimensional Span as a controlled technical headword whose equation layer, proof route, and registry role are interpreted together. \textbf{Derivable consequences:} The entry preserves the named object as a reusable definition, connects it to the local proof route, and records the consequences implied by the equation chain. \textbf{Lemma support:} Branch Neighborhood Declaration; Hyper-Dimensional Measure; Normalized Span.\par
\endgroup
% END TZP CONTEXTUAL MEANING

\tzpsection{Technical Interpretation}
This quantity measures inter-modal transfer across a hypersurface, scaled by the fundamental TZP loop unit.

\tzpsection{Equation Grounding}
Equation anchors: T sub IM ;= ; ratio 1 u sub T integral sub Sigma phi sub m sub 1 to the power * , L , phi sub m sub 2 , d Sigma, T sub IM ;= ; ratio 1 u sub T integral sub Sigma phi sub m sub 1 to the power * , L , phi sub m sub 2 , d Sigma,.

\tzpsection{Interpretive Role}
This quantity gives later actions and densities a corrected integration measure near the zero point. It is the bridge from local dimensional ambiguity to concrete higher-dimensional bookkeeping.
\endgroup

\tzpsection{TZP Thesaurus}
\begingroup\footnotesize\sloppy
\begin{enumerate}[leftmargin=1.35em,itemsep=1pt,topsep=1pt,parsep=0pt]
\item \textbf{Hyper Cross-Dimensional Span}: entry-specific alternate wording for indexing, related literature, and local formal context.
\item \textbf{Cross-Dimensional Formal Analogue}: entry-specific alternate wording for indexing, related literature, and local formal context.
\item \textbf{Cross-Dimensional Terminology}: entry-specific alternate wording for indexing, related literature, and local formal context.
\end{enumerate}
\endgroup

\tzpsection{Encyclopedia Mathematica}
\begingroup\footnotesize\sloppy
The visual plate below is reserved for the source-specific equation and progression graphic for [ID 0009]. It is included only as a companion to the canonical equation, not as a substitute for the formal text.
\par\endgroup
\ifTZPDarkEdition
\tzpdictionaryvisualband{D:/TZPID/kdp_workflow/build/tikz_figures/ID0009_dictionary_figure_dark.pdf}
\fi

\clearpage
\begingroup\tzpsupportsetup
\noindent\textbf{\fontsize{10.8pt}{12.8pt}\selectfont [ID 0009] Constructive Logic and Proof Certificate}\par
\noindent\textbf{\fontsize{10pt}{12pt}\selectfont Hyper-Dimensional Span}\par
\vspace{0.45em}
\begingroup
\footnotesize
\setlength{\parskip}{0.22em}
\setlength{\parindent}{0pt}
\noindent\textbf{Curry--Howard--Lambek Interpretation}\par
Under the Curry--Howard--Lambek correspondence, this registry entry is read as a typed proof
object: the stated source condition supplies the proposition, the derivation supplies the
morphism, and the proof certificate records the machine handles needed to check the obligation
in the formal registry.

\vspace{0.45em}
\noindent\textbf{CHL Proof}\par
\textbf{Statement:} the hyper-dimensional span measures the effective volume available in the neighborhood
where local dimension is allowed to branch around the zero p...

\textbf{Logic Validation:}\\
\textbf{Proposition:} ``Hyper-Dimensional Span is valid as a tensorial/geometric statement when the stated
tensor fields are well formed on the specified domain.''\\
\textbf{Proof:} The source statement gives a metric, curvature, or tensor relation. The CHL reading
validates it by checking that the tensorial objects have compatible domains, indices,
and transformation behavior.

\textbf{VERDICT:} \ensuremath{\checkmark}\ \textbf{TRUE} --- tensorial well-formedness validation

\textbf{Formal Status:} \ensuremath{\checkmark}\ THEORETICAL -- source-backed CHL proof obligation

\textbf{Physical Status:} THEORETICAL -- requires model-specific empirical interpretation
\par\vspace{0.45em}
\noindent{\tiny\textbf{Isabelle/HOL.} \tzpplaincode{physics\_validity\_from\_model, external\_validity\_package, registry\_number\_matches, equation\_count\_matches}}\par
\ifTZPDarkEdition
\tzpproofvisualband{D:/TZPID/kdp_workflow/build/tikz_figures/ID0009_proof_certificate_figure_dark.pdf}
\fi
\endgroup
\endgroup
